\begin{examplebox}

	\textbf{\textcolor{CExample}{例 1（基础题）}}

	\textbf{\textcolor{CExample}{题目：}}
	已知椭圆
	\[
		\frac{x^{2}}{25}+\frac{y^{2}}{9}=1
	\]
	绕 $x$ 轴旋转，求旋转椭球的表面积。

	\medskip

	\textbf{\textcolor{CExample}{解题步骤：}}

	\begin{description}[style=nextline, leftmargin=2.8em, labelsep=0.5em, itemsep=0.45em]

		\item[\textbf{第一步（读取参数）：}]
		      由椭圆方程得
		      \[
			      a=5,\qquad b=3.
		      \]

		      因为 $a>b$，所以这是长椭球情形。

		\item[\textbf{第二步（选择公式）：}]
		      当 $a>b$ 时，使用公式
		      \[
			      S
			      =
			      2\pi b^{2}
			      +
			      \frac{2\pi ab}{e}\arcsin e,
			      \qquad
			      e=\sqrt{1-\frac{b^{2}}{a^{2}}}.
		      \]

		\item[\textbf{第三步（计算离心率）：}]
		      \[
			      e
			      =
			      \sqrt{1-\frac{b^{2}}{a^{2}}}
			      =
			      \sqrt{1-\frac{9}{25}}
			      =
			      \sqrt{\frac{16}{25}}
			      =
			      \frac45.
		      \]

		\item[\textbf{第四步（代入公式）：}]
		      \[
			      \begin{aligned}
				      S
				       & =
				      2\pi\cdot 9
				      +
				      \frac{2\pi\cdot 5\cdot 3}{4/5}
				      \arcsin\frac45
				      \\
				       & =
				      18\pi
				      +
				      \frac{30\pi}{4/5}
				      \arcsin\frac45
				      \\
				       & =
				      18\pi
				      +
				      \frac{75\pi}{2}
				      \arcsin\frac45.
			      \end{aligned}
		      \]

	\end{description}

	\textbf{\textcolor{CExample}{答案：}}
	\[
		\boxed{
			S
			=
			18\pi
			+
			\frac{75\pi}{2}
			\arcsin\frac45
		}.
	\]

	\medskip

	\textbf{\textcolor{CExample}{例 2（稍变形）}}

	\textbf{\textcolor{CExample}{题目：}}
	已知椭圆
	\[
		\frac{x^{2}}{9}+\frac{y^{2}}{25}=1
	\]
	绕 $x$ 轴旋转，求旋转椭球的表面积。

	\medskip

	\textbf{\textcolor{CExample}{解题步骤：}}

	\begin{description}[style=nextline, leftmargin=2.8em, labelsep=0.5em, itemsep=0.45em]

		\item[\textbf{第一步（读取参数）：}]
		      由椭圆方程得
		      \[
			      a=3,\qquad b=5.
		      \]

		      因为 $a<b$，所以这是扁椭球情形。

		\item[\textbf{第二步（选择公式）：}]
		      当 $a<b$ 时，使用公式
		      \[
			      S
			      =
			      2\pi b^{2}
			      +
			      \frac{2\pi a^{2}}{e}\operatorname{artanh}e,
			      \qquad
			      e=\sqrt{1-\frac{a^{2}}{b^{2}}}.
		      \]

		\item[\textbf{第三步（计算离心率）：}]
		      \[
			      e
			      =
			      \sqrt{1-\frac{a^{2}}{b^{2}}}
			      =
			      \sqrt{1-\frac{9}{25}}
			      =
			      \frac45.
		      \]

		\item[\textbf{第四步（计算反双曲正切）：}]
		      \[
			      \operatorname{artanh}e
			      =
			      \frac12\ln\frac{1+e}{1-e}.
		      \]

		      代入 $e=\frac45$，得到
		      \[
			      \operatorname{artanh}\frac45
			      =
			      \frac12\ln\frac{1+\frac45}{1-\frac45}
			      =
			      \frac12\ln 9
			      =
			      \ln 3.
		      \]

		\item[\textbf{第五步（代入公式）：}]
		      \[
			      \begin{aligned}
				      S
				       & =
				      2\pi\cdot 25
				      +
				      \frac{2\pi\cdot 9}{4/5}\ln 3
				      \\
				       & =
				      50\pi
				      +
				      \frac{18\pi}{4/5}\ln 3
				      \\
				       & =
				      50\pi
				      +
				      \frac{45\pi}{2}\ln 3.
			      \end{aligned}
		      \]

	\end{description}

	\textbf{\textcolor{CExample}{答案：}}
	\[
		\boxed{
			S
			=
			50\pi
			+
			\frac{45\pi}{2}\ln 3
		}.
	\]

	\medskip

	\textbf{\textcolor{CExample}{例 3（综合应用）}}

	\textbf{\textcolor{CExample}{题目：}}
	某小行星的截面椭圆方程为
	\[
		\frac{x^{2}}{4}+\frac{y^{2}}{9}=1
	\]
	它绕 $x$ 轴旋转形成旋转椭球。求该小行星表面积的表达式。

	\medskip

	\textbf{\textcolor{CExample}{解题步骤：}}

	\begin{description}[style=nextline, leftmargin=2.8em, labelsep=0.5em, itemsep=0.45em]

		\item[\textbf{第一步（读取参数）：}]
		      由椭圆方程得
		      \[
			      a=2,\qquad b=3.
		      \]

		\item[\textbf{第二步（判断类型）：}]
		      因为 $a<b$，所以该旋转椭球为扁椭球。

		      使用公式
		      \[
			      S
			      =
			      2\pi b^{2}
			      +
			      \frac{2\pi a^{2}}{e}\operatorname{artanh}e.
		      \]

		\item[\textbf{第三步（计算离心率）：}]
		      \[
			      e
			      =
			      \sqrt{1-\frac{a^{2}}{b^{2}}}
			      =
			      \sqrt{1-\frac49}
			      =
			      \sqrt{\frac59}
			      =
			      \frac{\sqrt5}{3}.
		      \]

		\item[\textbf{第四步（代入公式）：}]
		      \[
			      \begin{aligned}
				      S
				       & =
				      2\pi\cdot 9
				      +
				      \frac{2\pi\cdot 4}{\sqrt5/3}
				      \operatorname{artanh}\frac{\sqrt5}{3}
				      \\
				       & =
				      18\pi
				      +
				      \frac{8\pi}{\sqrt5/3}
				      \operatorname{artanh}\frac{\sqrt5}{3}
				      \\
				       & =
				      18\pi
				      +
				      \frac{24\pi}{\sqrt5}
				      \operatorname{artanh}\frac{\sqrt5}{3}.
			      \end{aligned}
		      \]

	\end{description}

	\textbf{\textcolor{CExample}{答案：}}
	\[
		\boxed{
			S
			=
			18\pi
			+
			\frac{24\pi}{\sqrt5}
			\operatorname{artanh}\frac{\sqrt5}{3}
		}.
	\]

\end{examplebox}